\documentclass{handout}

\assignment{Homework 6}
\title{Standing and travelling waves}
\duedate{Wednesday October 22, 2025}

\begin{document}
\maketitle

\hwinstructions

\begin{questions}
    \question \textbf{Explaining standing waves:} Why does wave speed affect standing waves? They don't seem to be moving, so why would they have a speed?
    \begin{parts}
        \part Explain using the idea of function transformations why $\cos(kx)\cos(\omega t)$ represents a standing wave.
        \part Sinusoids follow a special identity: $\cos(a)\cos(b)=\frac{1}{2}(\cos(a-b) + \cos(a+b))$. Use this identity to separate a standing wave into two traveling waves.
        \part If the two traveling waves obey the relationship $c=\omega/k$, explain why the standing wave must obey this relationship as well.
        \part Which plays the more fundamental role in determining which waves can be standing for a particular cavity, $k$ or $\omega$?
    \end{parts}

    \question \textbf{Standing waves:} A string transmits waves at a speed of 264 m/s, is 30 cm long, and is fixed at both ends.
    \begin{parts}
        \part What is the longest wavelength standing wave the string can support?
        \part What will be the frequency of the fundamental mode (the dominant tone if you pluck the string)? How did you determine this number?
        \part What will be the frequency of the sound that reaches your ears?
        \part The speed of sound in air is 343 m/s. What is the wavelength of the sound waves that reach your ears?
        \part Where could you place your finger along this string to make the third harmonic? Explain your answer. What would be the frequency of this wave?
    \end{parts}

    \question \textbf{Closed/open pipe harmonics:} When you blow across the open end of a pipe, you randomly excite standing pressure waves in the pipe.
    \begin{parts}
        \part In a pipe that is \textbf{open} on both ends, the pressure must be the same as the surrounding atmosphere at the ends of the pipe. What are the first three fundamental wavelengths and frequencies for a pipe in terms of its length $L$ and the speed of sound in air $c$? Draw these standing waves.
        \part Recall that Bernoulli's principle states that the velocity of a particle of fluid is proportional to how fast the pressure is changing at that location. Draw the standing waves of \emph{velocity} corresponding to the pressure waves drawn above.
        \part What is the direction of the velocities discussed in the previous part?
        \part For a pipe that is \textbf{closed} on one end but \textbf{open} on the other, the \textbf{velocity} of particles at the open end must be zero. Draw the first three standing waves of \emph{velocity} in such a pipe.
        \part For the open/closed pipe standing waves drawn above, draw the corresponding pressure waves.
    \end{parts}
\end{questions}

\end{document}
